\setcounter{chapter}{1}
\chapter{Architecture and Design}
\minitoc
\graphicspath{{Chapter2/figures/}}

%==============================================================================
\pagestyle{fancy}
\fancyhf{}
\fancyhead[R]{\bfseries\rightmark}
\fancyfoot[R]{\thepage}
\renewcommand{\headrulewidth}{0.5pt}
\renewcommand{\footrulewidth}{0pt}
\renewcommand{\chaptermark}[1]{\markboth{{\chaptername~\thechapter.\ #1}}{}}
\renewcommand{\sectionmark}[1]{\markright{\thechapter.\thesection~ #1}}

\begin{spacing}{1.2}
%==============================================================================

\section*{Introduction}
% TODO:
%   - Frame this chapter as the architectural narrative that the release
%     chapters (III to VIII) will reference back to.
%   - State the three forces shaping the architecture: integration-first
%     (no document migration), compliance-by-construction (immutable audit
%     trail), and single-tenant-per-customer (DACH Mittelstand sensibility).
%   - Preview the sections: drivers, CQRS/ES rationale, comby framework,
%     domain model, system topology, security & compliance, i18n.

%--------------------------------------------------------------------
\section{Architectural drivers}
% TODO: Recap the NFRs from Chapter I as ARCHITECTURAL DRIVERS, one
% paragraph each. These are the constraints that any solution must
% satisfy regardless of release scope.
%   - Auditability and immutability (GoBD, ISO 27001).
%   - Integration-first (documents stay in Drive/SharePoint -- non-negotiable).
%   - EU data residency (GDPR; the entire stack hosts in Germany).
%   - Single-tenant economics (DACH Mittelstand customers expect a
%     dedicated logical tenant, not a shared multi-tenant pool).
%   - Maintainability under solo development.
% These drivers will be referenced explicitly in every Phase chapter as
% the justification for each decision taken at delivery time.
\dots

%--------------------------------------------------------------------
\section{Why CQRS and Event Sourcing for CLM?}
% TODO: Architectural reasoning, not implementation. Cite Young 2010,
% Fowler CQRS/ES, Vernon 2013, Evans 2003.
%   - The CLM domain is naturally event-driven: ``contract signed'',
%     ``approval granted'', ``deadline crossed'' -- these are events,
%     not states.
%   - Event Sourcing yields the audit trail for free; it is not a
%     feature bolted on top, it is the storage model.
%   - Command/Query separation maps cleanly to the read-heavy nature
%     of dashboards and the write-careful nature of approvals.
%   - Trade-offs honestly stated: eventual consistency in read models,
%     event-schema versioning, learning curve.
%   - Why not alternatives: CRUD with audit columns (audit trail is a
%     bolt-on, not a guarantee), document databases (no transactional
%     model for approvals), traditional workflow engines (vendor lock-in).
\dots

%--------------------------------------------------------------------
\section{The comby framework}
% TODO: Engineer-level presentation. What comby IS, what role it plays,
% why it was chosen.
%   - One-paragraph definition: an in-house Go framework providing the
%     command/query buses, the event store abstraction, the projection
%     runner, and the reactor scaffolding for CQRS+ES applications at
%     Gradient Zero.
%   - Position in the build vs. buy axis: not a developer convenience
%     library but a strategic architectural asset of the company.
%   - What CLMPilot inherits from comby and what it builds on top.
%   - Alternative considered: EventStoreDB + custom Go code; Axon
%     (Java only); EventFlow (.NET); bespoke. Justification for comby
%     in one paragraph.

\subsection{Role in the architecture}
\dots

\subsection{Justification of choice}
% Inline comparison table -- the trade-off discussion is the engineer's
% deliverable, not the API surface.
\dots

%--------------------------------------------------------------------
\section{Domain model overview}
% TODO: Show the four bounded contexts at the conceptual level only.
% NO aggregate-by-aggregate event lists or command lists here -- those
% live in Chapter IV.
%   - Diagram: the four aggregates and their relationships
%     (Contract, Approval, Deadline, DocumentLink), plus the boundary
%     between them.
%   - One paragraph per aggregate, focusing on the BUSINESS responsibility,
%     not the event types: Contract = the central record; Approval =
%     the workflow gate; Deadline = the temporal awareness; DocumentLink
%     = the integration boundary.
%   - The cross-cutting concerns (identity, audit log) live outside
%     the four aggregates as platform-level concerns.

\subsection{Bounded contexts}
\begin{figure}[!ht]\centering
  % \includegraphics[scale=0.7]{domain-model-overview.png}
  \caption{The four bounded contexts of CLMPilot and their relationships.}
  \label{fig:domains}
\end{figure}

\subsection{Aggregate responsibilities}
\dots

%--------------------------------------------------------------------
\section{System topology}
% TODO: One diagram, brief commentary. The topology is the answer to
% ``where does code run and where does data live''.
%   - Four repositories, four services: backend (Go), frontend (Vue),
%     intelligence (Python stub for MVP), deployment (compose stack).
%   - Two stateful stores: PostgreSQL event store, Redis (sessions
%     + MFA tokens).
%   - One reverse proxy (nginx on the host).
%   - External: SendGrid for email, Google Drive / SharePoint as
%     integration targets, eventually DocuSign (Phase 2) and
%     ERP systems (Phase 3).

\subsection{Deployment topology}
\begin{figure}[!ht]\centering
  % \includegraphics[scale=0.7]{system-topology.png}
  \caption{CLMPilot deployment topology and external integrations.}
  \label{fig:topology}
\end{figure}

\subsection{Why a 4-repository split}
% Engineer-level reasoning: separation of release cadence, separation of
% concerns (backend / UI / intelligence / ops), team-of-one autonomy.
\dots

%--------------------------------------------------------------------
\section{Security and compliance model}
% TODO: Authentication and authorization as ARCHITECTURE, not
% as ``how MFA was wired in''.
%   - Identity model: tenant, user, role, permission.
%   - Authentication factors: password + email MFA OTP (SendGrid HTTP API).
%   - Authorization: RBAC scoped per tenant.
%   - Compliance positioning: GoBD compliance comes from the immutable
%     event log (Section 2 above); GDPR from EU hosting + per-tenant
%     data isolation.
%   - Audit trail availability is a CONSEQUENCE of ES, not an additional
%     feature -- this is the central architectural sales pitch.

\subsection{Identity and access model}
\dots

\subsection{Compliance positioning}
\dots

%--------------------------------------------------------------------
\section{Internationalisation strategy}
% TODO: Architectural decision, not i18n file format.
%   - EN as primary language; DE and FR scaffolded for DACH expansion.
%   - Server-side message catalogue (used by email + audit log) and
%     client-side message catalogue (used by the SPA) kept in sync.
%   - User-locale resolution: explicit setting first, browser hint
%     second, EN default.
\dots

%--------------------------------------------------------------------
\section*{Conclusion}
% TODO: Recap the architectural foundation. Bridge to Chapter III:
% ``With the architecture set, the next chapters walk through the
% release-based delivery. Chapter III covers the engineering
% foundations established in Phase 0 before any feature code was
% written.''
\dots

%==============================================================================
\end{spacing}
